\section{模型近似的五大缺陷}

\subsection{语义定义的根本困境}

张钹院士提出了一个深刻的观点：当前大语言模型中出现的很多怪现象，\textbf{不能简单归结于数据问题}。

\begin{warningbox}{被广泛误解的归因}
``很多人发现机器里出现很多怪现象，都归结于数据问题。这是完全错误的。现在很多东西是由于\textbf{模型近似}引起的。''世界上有七八种不同哲学学派对``语义''的定义，没有任何一种是科学上完备的。分布式语义只是其中一种近似。
\end{warningbox}

\subsection{五大缺陷}

由于语义定义的不完备性，大语言模型必然存在以下五个缺陷：

\begin{importantbox}{大语言模型的五大固有缺陷}
\begin{enumerate}
\item \textbf{指称缺陷}：模型不能真正指向现实世界中的实体
\item \textbf{真实和因果缺陷}：模型缺乏因果推理能力，只能捕捉统计相关性
\item \textbf{语用缺陷}：模型无法真正理解语言在特定社会情境中的使用规则
\item \textbf{多义和动态缺陷}：词义的多样性和动态变化超出了静态向量表示的能力
\item \textbf{语境和闭环行为缺陷}：模型缺乏与物理环境的闭环交互能力
\end{enumerate}
\end{importantbox}

张钹院士强调：这五个缺陷是\textbf{固有的、必然存在的}——不是通过更多数据或更大模型就能消除的，因为问题的根源在于语义定义本身的不完备。


